\begin{theorem}\label{dedekind-30} The identical transformation of a
system is always a similar transformation.

\end{theorem}
\begin{theorem}\label{dedekind-31} If $\phi$ is a similar transformation of
$S$ and $\psi$ a similar transformation of $\phi(S)$, then is the
transformation $\psi\phi$ of $S$, compounded of $\phi$ and $\psi$, a similar
transformation, and the associated inverse transformation
$\overline{\psi}\,\overline{\phi} = \overline{\phi}\,\overline{\psi}$.
\end{theorem}
\begin{proof}For to different elements $a$, $b$ of $S$ correspond
different transforms $a' = \phi(a)$, $b'=\phi(b)$, and
to these again different transforms $\psi(a')=\psi\phi(a)$,
$\psi(b') =\psi\phi(b)$ and therefore $\psi\phi$ is a similar transformation.
Besides every element $\psi\phi(s) = \psi(s')$ of the
system $\psi\phi(S)$ is transformed by $\overline{\psi}$ into $s' = \phi(s)$ and
% [File: 060.png]
this by $\overline{\phi}$ into $s$, therefore $\psi\phi(s)$ is transformed by
$\overline{\phi}\,\overline{\psi}$ into $s$, which was to be proved.

\end{proof}
